\startprefacepage
%% Макрос для введения

Лазерная сварка представляет собой способ неразъёмного соединения материалов, при котором шов формируется сфокусированным лазерным излучением. По сравнению с традиционными электродуговыми методами она требует более сложного и дорогостоящего оборудования, однако это компенсируется локальностью теплового воздействия, малой зоной термического влияния и высокой точностью. Поэтому лазерная сварка особенно перспективна для автоматизированного производства, где важны повторяемость процесса и возможность управления без участия сварщика~\cite{wu2022progress,huang2025laser}.

Вместе с тем лазерная сварка остаётся сложным и чувствительным к режимам процессом. На качество соединения влияют мощность излучения, скорость перемещения пучка, положение фокуса, свойства материала, геометрия стыка и состояние поверхности. В зоне обработки одновременно протекают нагрев, плавление, испарение, колебание парогазового канала, течение расплава, формирование плазменного факела и кристаллизация металла. Поэтому связь между параметрами процесса и качеством шва носит нелинейный характер, а небольшие отклонения могут приводить к непровару, пористости и другим дефектам~\cite{huang2025laser,wu2022progress}.

Следовательно, для реализации преимуществ лазерной сварки недостаточно заранее задать режимы и воспроизвести траекторию движения лазерной головки. Необходима система, способная контролировать процесс во время сварки и при необходимости корректировать параметры воздействия. В этом случае высокая управляемость лазерного источника становится основой системы замкнутого цикла, где качество соединения поддерживается обратной связью по текущему состоянию зоны сварки~\cite{konuk2011closed,zhang2018adaptive}.

Одним из наиболее перспективных источников такой обратной связи являются оптические сигналы из зоны сварки. Излучение расплава, плазмы, области парогазового канала и отражённое лазерное излучение несут информацию о состоянии процесса, глубине проплавления и возникновении дефектов. Поэтому фотодиоды, спектрометры и другие оптические датчики могут использоваться как основа для оценки качества и последующего управления лазерной сваркой в реальном времени~\cite{konuk2011closed,huang2025laser}.

\endinput